\documentclass[10pt,compress,t,notes=noshow,xcolor=table]{beamer}

\input{../../style/preamble}
% Defines macros and environments

\title{Interpretable Machine Learning}
% \author{LMU}
%\institute{\href{https://compstat-lmu.github.io/lecture_iml/}{compstat-lmu.github.io/lecture\_iml}}

\begin{document}

\titlemeta{
Interpretable Models 1
}{
Inherently Interpretable Models - Motivation
}{
figure/whitebox
}{
%\item What characteristics does an interpretable model have?
\item Why should we use interpretable models?
\item Advantages and disadvantages of interpretable models
}

\begin{framev}[align=top]{Motivation}
%\bigskip
\splitVTT[0.58]{
\spacer[0.6]
\begin{itemizeM}
\item Achieving interpretability by using interpretable models is the most straightforward approach
\item Classes of models deemed interpretable:
\begin{itemizeS}[small]
\item (Generalized) linear models (LM, GLM)
\item Generalized additive models (GAM)
\item Decision trees
\item Rule-based learning
\item Model-based / component-wise boosting
\item ...
\end{itemizeS}
\end{itemizeM}
}{
\spacer[0]
\imageC[0.8]{figure/main_effect_lm_temp.pdf}
\centering
$\leadsto$ LM provides straightforward interpretation
\par
}
\begin{itemizeM}
\item Often there is a trade-off between interpretability and model performance
%(but not always)
\end{itemizeM}
\imageC[0.55]{figure/performance_vs_interpretability.pdf}
% Quelle: https://docs.google.com/presentation/d/12ZPrTjBKEUT-7drdyUJCQGK0oHVDtYIVd2_6byE62f0/edit?usp=sharing
\end{framev}


\begin{framev}[align=top]{Advantages}
\splitVTT[0.78]{
\spacer[0.6]
\begin{itemizeM}
\item<+-> Interpretable models are transparent by design, making many model-agnostic explanation methods unnecessary\\
$\leadsto$ Eliminates an extra source of estimation error
\item<+-> They often have few hyperparameters and are structurally simple (e.g., linear, additive, sparse, monotonic)\\
$\leadsto$ Easy to train, fast to tune, straightforward to explain
% \item Some interpretable models estimate monotonic effects \\
% %the monotonicity constraint \\
% $\leadsto$ Simple to explain as larger feature values always lead to higher (or smaller) outcomes (e.g., GLMs)
\item<+-> Many people are familiar with simple interpretable models\\
$\leadsto$ Increases trust, facilitates communication of results
%\item Implementations are available in many programming languages \\
%$\leadsto$ Simple models are easier to deploy in practice or implement from scratch
\end{itemizeM}
}{
\spacer[0]
\begin{center}
\begin{tikzpicture}[scale=0.6, transform shape]
\usetikzlibrary{arrows}
\usetikzlibrary{shapes}
\tikzset{treenode/.style={draw, circle, font=\small}}
\tikzset{line/.style={draw, thick}}
\node [treenode, draw=red] (a0) {$a_0$};
\node [treenode, below=0.75cm of a0, xshift=-1cm] (a1) {$a_1$};
\node [treenode, draw=red, below=0.75cm of a0, xshift=1cm] (a2) {$a_2$};
\node [treenode, draw=red, below=0.75cm of a2, xshift=-1cm] (a3) {$a_3$};
\node [treenode, below=0.75cm of a2, xshift=1cm] (a4) {$a_4$};
\node [treenode, below=0.75cm of a3, xshift=-1cm] (a5) {$a_5$};
\node [treenode, below=0.75cm of a3, xshift=1cm] (a6) {$a_6$};
\path [line] (a0.south) -- +(0,-0.4cm) -| (a1.north) node [midway, above] {$x_1<0.3$};
\path [line] (a0.south) -- +(0,-0.4cm) -| (a2.north) node [midway, above] {$x_1\geq0.3$};
\path [line] (a2.south) -- +(0,-0.4cm) -| (a3.north) node [midway, above] {$x_1<0.6$};
\path [line] (a2.south) -- +(0,-0.4cm) -| (a4.north) node [midway, above] {$x_1\geq0.6$};
\path [line] (a3.south) -- +(0,-0.4cm) -| (a5.north) node [midway, above] {$x_2<0.2$};
\path [line] (a3.south) -- +(0,-0.4cm) -| (a6.north) node [midway, above] {$x_2\geq0.2$};
\end{tikzpicture}
\end{center}
\spacer[0.25]
\only<2->{
\centering
\image[0.95]{figure/main_effect_lm_temp.pdf}
\par
}
}
\end{framev}


\begin{framev}[align=top]{Disadvantages \& Limitations}
\splitVTT[0.8]{
\spacer[0.6]
\begin{itemizeM}
\item<1-> Often require assumptions about data / model structure\\%Often certain assumptions about data / model structure required\\
$\leadsto$ If assumptions are wrong, models may perform bad
%\item In high-dimensional settings, it may not be feasible to find and define all feature relationships
%\item If the wrong assumptions are made, interpretable models may have a bad predictive performance
%\pause
\item<2-> Interpretable models may also be hard to interpret, e.g.:
\begin{itemizeS}
\item LM with lots of features and interactions
\item Decision trees with huge tree depth
\end{itemizeS}
\end{itemizeM}
}{
\spacer[0]
\imageC{figure/lm_bad_fit.pdf}
\spacer[0]
}
\spacer[0.25]
\begin{itemizeM}
\item<3->
Often do not automatically model complex relationships due to limited flexibility % didn't find a way to fig the hanging word
\\
e.g., high-order main or interaction effects need to be specified manually in LM
%\pause
\item<4-> Inherently interpretable models do not address all explanation needs\\
$\leadsto$ Complementary model-agnostic methods (e.g., counterfactuals) remain valuable for specific tasks
\end{itemizeM}
\end{framev}


\begin{framev}[align=top]{Further Comments}
\begin{itemizeM}
%\itemsep1em
\item<1-> Some researchers advocate for inherently interpretable models instead of explaining black boxes after training \furtherreading{RUDIN2019}
\begin{itemizeS}
\item Built-in interpretation\\
$\leadsto$ fewer risks from misleading post-hoc explanations
\item Good performance possible with effort on preprocessing and/or feature engineering
\item But interpretability depends on meaning of created features\\
$\leadsto$ E.g., PCA keeps models linear, but yields hard-to-interpret components
%, or data cleaning
\end{itemizeS}
%\pause
%\footnote[frame]{Rudin, C. Stop explaining black box machine learning models for high stakes decisions and use interpretable models instead. Nat Mach Intell 1, 206–215 (2019).}
% {https://www.nature.com/articles/s42256-019-0048-x}
\item<2-> Limitation: Less suited for complex data requiring end-to-end learning
\begin{itemizeS}
\item Applies to image, text, or sensor data where features must be learned from raw input
\item Manual extraction of interpretable features is difficult\\
$\Rightarrow$ Information loss and lower performance
\end{itemizeS}
%(e.g., for image / text data extraction of good features is hard $\leadsto$ information loss = bad performance)
%(e.g., image and text data would require good feature extraction steps $\leadsto$ information loss)
%\pause
\end{itemizeM}
\end{framev}


\begin{framev}[align=top]{Recommendation}
\begin{itemizeM}
\item Begin with the simplest model appropriate for the task
\item Increase complexity only if necessary to meet performance requirements\\
$\leadsto$ Typically reduces interpretability and requires model-agnostic explanations
\item Choose the simplest model with sufficient accuracy\\
$\leadsto$ Occam's razor
\end{itemizeM}
\spacer[0.25]
% \begin{columns}[T, totalwidth=\textwidth]
% \begin{column}{0.5\textwidth}
% \centering \textbf{Bike Data}
%     \begin{table}[ht]
%         \centering
%         \begin{tabular}{lrr}
%             \hline
%             Model & RMSE & $R^2$ \\
%             \hline
%             LM & 142.39 & 0.38 \\
%             Tree & 98.71 & 0.70 \\
%             Random Forest & 57.07 & 0.90 \\
%             Boosting & 197.42 & -0.18 \\
%             \hline
%         \end{tabular}
%     \end{table}
% \end{column}
% \begin{column}{0.5\textwidth}
% \centering \textbf{Boston Housing Data}
%     \begin{table}[ht]
%         \centering
%         \begin{tabular}{lrr}
%             \hline
%             Model & RMSE & $R^2$ \\
%             \hline
%             LM & 0.43 & 1.00 \\
%             Tree & 1.81 & 0.96 \\
%             Random Forest & 1.77 & 0.96 \\
%             Boosting & 16.89 & -2.43 \\
%             \hline
%         \end{tabular}
%     \end{table}
% \end{column}
% \end{columns}
% code for tables in rsrc/tab_ml_comparison.R
\centering
\textbf{Bike Data, 4-fold CV}\\
\spacer[0.25]
\begin{table}
\centering
\begin{tabular}{lrr}
\hline
Model & RMSE & $R^2$ \\
\hline
LM & 800.15 & 0.83 \\
Tree & 981.83 & 0.74 \\
Random Forest & 653.25 & 0.88 \\
Boosting (tuned) & 638.42 & 0.89 \\
\hline
\end{tabular}
\end{table}
\end{framev}

\endlecture
\end{document}
